%% Lecture d'une affixe dans le plan complexe : les cercles concentriques donnent le
%% module (distance OM), l'arc orienté depuis l'axe réel donne l'argument.
%% Volontairement, aucune valeur n'est écrite : c'est à l'élève de lire M.
\documentclass[tikz,border=4pt]{standalone}
\usepackage{liaisons}
\begin{document}
\begin{tikzpicture}[x=0.52cm,y=0.52cm,
  axe/.style={-{Stealth[length=5pt]}, line width=0.6pt},
  rappel/.style={line width=0.6pt, dash pattern=on 3pt off 2pt, black!55},
  grad/.style={font=\scriptsize, inner sep=0pt}]

%% ---------------- cinq cercles concentriques -------------------------------
\foreach \r in {1,2,3,4,5} { \draw[black!35, line width=0.5pt] (0,0) circle (\r); }

%% ---------------- axes -------------------------------------------------------
\draw[axe] (-5.6,0) -- (6.1,0);
\draw[axe] (0,-5.6) -- (0,6.1);
\node[font=\small, anchor=south east, inner sep=3pt] at (6.15,0.1) {axe réel};
\node[font=\small, anchor=east, align=right, inner sep=3pt, fill=white] at (-0.2,5.1)
     {axe des\\imaginaires purs};
\node[font=\small, anchor=north east, inner sep=3pt] at (-0.1,-0.1) {$O$};
\foreach \r in {1,2,3,4,5} {
  \draw[line width=0.5pt] (\r,0) -- ++(0,-0.07cm);
  \node[grad, anchor=north, yshift=-3pt, fill=white, inner sep=1pt] at (\r,0) {$\r$};}

%% ---------------- le point M (sur le cercle de rayon 5, quatrième quadrant) --
\draw[rappel] (3.5355,-3.5355) -- (3.5355,0);
\draw[rappel] (3.5355,-3.5355) -- (0,-3.5355);
\draw[solideA, line width=1.5pt] (0,0) -- (3.5355,-3.5355);
\fill (3.5355,-3.5355) circle (2.4pt);
\node[font=\small, anchor=south west, inner sep=3pt] at (3.5355,-3.5355) {$M$};

%% ---------------- arc orienté de l'axe réel vers [OM] -----------------------
\draw[solideE, line width=1.2pt, -{Stealth[length=6pt]}] (2.3,0) arc (0:-39:2.3);
\end{tikzpicture}
\end{document}
